
\def\PUDcodigo{EME112}

\def\PUDcodacad{04506.29}

\def\PUDdisciplina{Soldagem}

\def\PUDsemestre{S6}

\def\PUDhoras{80}

%NOVOS ---------------------------------------------------
\def\PUDhorasTeoria{50}
\def\PUDhorasPratica{30}

\def\PUDinterECA{---}

\def\PUDatuaisECA{---}

\def\PUDrecursos{Projetor multimídia; Quadro branco e pincel; Aulas em laboratório.}

%----------------------------------------------------------

\def\PUDcreditos{4}

\def\PUDprerequisito{ \hyperref[PUD:EME108]{EME108} }

\def\PUDpreacad{ \hyperref[PUD:EME108]{04506.20} }

\def\PUDobjetivos{Compreender a importância do processo de soldagem na manutenção industrial;  Conhecer os diferentes tipos de processos de soldagem a arco voltaico;  Compreender os equipamentos empregados nos diversos processos de soldagem a arco voltaico.  Avaliar os efeitos dos parâmetros de soldagem empregados na união de materiais.  Estar apto a realizar análises das soldagens empregadas na manutenção de equipamentos;  Fazer
a correta seleção do processo de soldagem para uma dada aplicação.}

\def\PUDementa{Princípios de Soldagem;  Terminologia;  Segurança na Soldagem;  Arco Elétrico;  Eletrodo Revestido;  TIG;  Transferência Metálica;  MIG/MAG;  Arame Tubular;  Arco Submerso;  Plasma;  Soldagem de Revestimento;  Metalurgia da Soldagem.}

\def\PUDconteudo{ & Princípios de soldagem. 
\\ 
 & Terminologia. 
\\ 
 & Arco elétrico. 
\\ 
 & Eletrodo revestido. 
\\ 
 & TIG. 
\\ 
 & Transferência metálica. 
\\ 
 & MIG/MAG. 
\\ 
 & Arame tubular. 
\\ 
 & Arco submerso. 
\\ 
 & Plasma. 
\\ 
 & Soldagem de revestimento. 
\\ 
 & Metalurgia da soldagem. }

\def\PUDmetodo{Aulas expositórias que podem ser teóricas e/ou práticas, onde as práticas no laboratório serão marcadas no decorrer da disciplina.}

\def\PUDavalia{A avaliação será feita com aplicação de provas teóricas e/ou práticas, além da possibilidade de inclusão de trabalhos e seminários no decorrer da disciplina.}

\def\PUDbibbasica{ & \href{www.biblioteca.ifce.edu.br/index.asp?codigo_sophia=24034}{Wainer}, E.; Brandi, S. D.; Décon, F.  Soldagem - Processos e Metalurgia.  1ª ed.  São Paulo, SP: Blucher, 1992.  494p.  ISBN: 9788521202387. 
\\ 
 & \href{www.biblioteca.ifce.edu.br/index.asp?codigo_sophia=41944}{Veiga}, E.  Soldagem de Manutenção.  1ª ed.  São Paulo, SP: Globus, 2011.  328p.  ISBN: 9788579810497. 
\\ 
 & \href{www.biblioteca.ifce.edu.br/index.asp?codigo_sophia=24061}{Chiaverini}, V.  Tecnologia Mecânica - Processos de Fabricação e Tratamentos.  vol.2.  2ª ed.  São Paulo, SP: Pearson Makron Books, 1986.  315p.  ISBN: 9780074500903. }

\def\PUDbibcomplementar{ & \href{www.biblioteca.ifce.edu.br/index.asp?codigo_sophia=24147}{Weiss}, Almiro.  Soldagem.  Curitiba, PR: Livro Técnico, 2010.  128p.  ISBN: 9788563687166.
%&  MARQUES, P.  V.;  MODENESI, P.  J.;  BRACARENSE, A.  Q.;  Soldagem - Fundamentos e Tecnologia.  3ª edição.  Belo Horizonte.  Editora UFMG.  362p.    ISBN 9788570417480. 
\\ 
 &  \href{www.biblioteca.ifce.edu.br/index.asp?codigo_sophia=61950}{Scotti}, Americo; Ponomarev, Vladimir.  Soldagem MIG/MAG: melhor entendimento, melhor desempenho.  2ª ed.  São Paulo, SP: Artliber, 2014. 288p.  ISBN: 9788588098428.       
\\ 
  & \href{www.biblioteca.ifce.edu.br/index.asp?codigo_sophia=39135}{Souza}, Sérgio Augusto de.  Composição química dos aços.  São Paulo, SP: Edgard Blücher, 1989.  134p. ISBN: 9788521203025. 
\\
 & \href{www.biblioteca.ifce.edu.br/index.asp?codigo_sophia=39995}{Chiaverini}, Vicente.  Aços e ferros fundidos : características gerais, tratamentos térmicos e principais tipos.  7º ed.  São Paulo, SP: ABM, 2012.  599p.  ISBN: 8586778486.
\\
 & \href{www.biblioteca.ifce.edu.br/index.asp?codigo_sophia=41211}{Padilha}, Angelo Fernando.  Encruamento, recristalização, crescimento de grão e textura.  3º ed.  São Paulo, SP: ABM, 2005.  232p.  ISBN: 858677880X. }
%&  QUITES, Almir M.;  Metalurgia na Soldagem dos Aços.  2ª Edição.  Florianópolis.  Editora SOLDASOFT.  2009.  304p.  ISBN : 9788589445054. }

\def\PUDautor{Rodrigo Freitas Guimarães}

\def\PUDdata{2013-06-19}

\def\PUDrevisao{1}

\def\PUDrevisor{Samuel Vieira Dias}

\def\PUDdatarevisao{2017-05-23}

\PUDcorpo
